% Generated draft: requires mathematical and source review.
\documentclass[11pt]{article}
\usepackage[margin=25mm]{geometry}
\usepackage{fontspec}
\setmainfont{Latin Modern Roman}
\usepackage{amsmath,amssymb,mathtools,unicode-math}
\setmathfont{Latin Modern Math}
\usepackage{enumitem,needspace,xurl,hyperref}
\setlength{\parindent}{0pt}
\setlength{\parskip}{5pt}
\setlength{\emergencystretch}{3em}
\newcommand{\N}{\mathbb{N}}
\newcommand{\Z}{\mathbb{Z}}
\newcommand{\Q}{\mathbb{Q}}
\newcommand{\R}{\mathbb{R}}
\newcommand{\C}{\mathbb{C}}
\newcommand{\F}{\mathbb{F}}
\newcommand{\A}{\mathbb{A}}
\newcommand{\PP}{\mathbb{P}}
\newcommand{\E}{\mathbb{E}}
\newcommand{\eps}{\varepsilon}
\newcommand{\dd}{\,\mathrm{d}}
\begin{document}
\section*{UnsolvedMath ID: 800306}
\textbf{Source problem code:} catalogue entry
\subsection{Definitions and mathematical statement}
Let \(S^1=\R/\Z\), and let \(\operatorname{Homeo}(S^1)\) denote the group of circle homeomorphisms. Let \(G\leq\operatorname{Homeo}(S^1)\) be finitely generated, and write \(\rho_0:G\hookrightarrow\operatorname{Homeo}(S^1)\) for its given action. An additional circle action of the abstract group \(G\) is a homomorphism \(\rho:G\to\operatorname{Homeo}(S^1)\). For a finite generating set \(S\subseteq G\), the usual \(C^0\) topology on the space of actions is the topology of uniform convergence of the homeomorphisms assigned to the generators and their inverses.

The supplied question asks whether the following implication holds: if \(\rho_0\) is \(C^0\) structurally stable and \(G\) admits infinitely many pairwise non-semiconjugate circle actions, must there exist \(a,b\in G\) such that
\[
\langle a,b\rangle\cong F_2,
\]
where \(F_2\) is the free group on two generators? Equivalently, the conclusion requires that every nonempty reduced word in \(a^{\pm1}\) and \(b^{\pm1}\) represent a nonidentity element of \(G\).

\emph{Unresolved qualifications.} The supplied record does not define \(C^0\) structural stability or semiconjugacy. In particular, it does not specify the equivalence required between nearby actions, the permitted regularity or monotonicity of a semiconjugating map, or whether semiconjugacy is intended as a symmetric equivalence relation. It also does not specify whether only orientation-preserving actions are allowed or whether the additional actions must be faithful. The implication above retains these source terms pending clarification; it is not a fully specified assertion until those conventions are fixed.
\subsection{Short English statement}
For a finitely generated group with a \(C^0\) structurally stable circle action, does the existence of infinitely many pairwise non-semiconjugate circle actions force a free subgroup of rank two? The source conventions for stability, semiconjugacy, orientation, and faithfulness require clarification. The supplied registry labels the problem open; that label is not independently verified.
\subsection{Sources}
\begin{itemize}
\item Local UnsolvedMath registry, supplied complete source record: permanent ID \(800306\), source problem code \texttt{catalogue entry}, title ``Group Actions on 1-Manifolds --- Question 23.'' The record attributes the question to Navas, with proposed year 2017, and supplies the status \texttt{open}.
\item Navas, \emph{Group Actions on 1-Manifolds: A List of Very Concrete Open Questions} (2017), source item recorded as ``Question 23 (\#22).'' Supplied source URL: \url{https://arxiv.org/abs/1712.06462}. The supplied provenance records access on 2026-10-16 and extraction method \texttt{arxiv\_latex}; these are supplied provenance claims, not independently verified access or source checking.
\end{itemize}
\end{document}
